\section{Host orchestration: one submission per chunk}
\label{sec:orch}

\subsection{The driver}

\code{OctreeTopologyCompute::computeResident} is the only entry point.  It
is called by \code{SceneRenderer::buildChunkResident} on a scene worker
thread.  The method holds a mutex for the whole chunk, so two workers never
share the staging buffers or the descriptor set.  The normal path is one
submission:

\begin{enumerate}
\item the host flattens the neighbourhood and derives the cross-edge
  capacities;
\item the provider reserves a vertex/index span in the packed pools;
\item the host uploads the node tables and records one geometry dispatch;
\item the host waits for the fence and reads the counters;
\item the host publishes the span, or grows it once and retries.
\end{enumerate}

\begin{figure}[htbp]
\centering
\begin{tikzpicture}[font=\small]
  \draw[->,thick] (0,0) -- (12.4,0) node[right]{time};
  \node[above] at (1.2,0) {CPU};
  \node[above] at (1.2,-1.2) {GPU};
  \draw[dashed] (1.2,-1.4)--(1.2,0.6);
  % CPU work
  \fill[cpu,draw=accent] (0.2,0.15) rectangle (2.6,0.65) node[midway,font=\scriptsize]{flatten + capacity};
  \fill[cpu,draw=accent] (2.8,0.15) rectangle (4.4,0.65) node[midway,font=\scriptsize]{reserve span};
  \fill[cpu,draw=accent] (4.6,0.15) rectangle (6.2,0.65) node[midway,font=\scriptsize]{upload};
  \fill[cpu,draw=accent] (8.4,0.15) rectangle (10.4,0.65) node[midway,font=\scriptsize]{read 10 counters};
  \fill[cpu,draw=accent] (10.6,0.15) rectangle (12.2,0.65) node[midway,font=\scriptsize]{publish / retry};
  % GPU work
  \fill[gpu,draw=accent2] (6.4,-1.0) rectangle (8.2,-0.5) node[midway,font=\scriptsize]{geometry pass};
  % arrows
  \draw[arrow] (4.4,0.4)--(4.6,0.4);
  \draw[arrow] (6.2,0.4)--(6.4,0.4)--(6.4,-0.4);
  \draw[arrow] (8.2,-0.75)--(8.4,-0.75)--(8.4,0.4);
  \draw[arrow] (10.4,0.4)--(10.6,0.4);
\end{tikzpicture}
\caption{One chunk's normal timeline.  The spans are reserved from the CPU
cross-edge bound before submission; the single geometry pass writes the
renderer's pools directly.  The host reads integers after the fence.}
\label{fig:timeline}
\end{figure}

\subsection{The input upload}

The input buffers are host-visible and coherent.  \code{uploadInput} writes:
\begin{itemize}
\item 128 bytes per node for the node table;
\item 4 bytes per cell for the scan list;
\item ten counters (indices, extent, overflow, drops, diagnostics, and the
  distinct vertex write count);
\end{itemize}
and resets the counters.  The upload is not a single \code{memcpy}: the
node record is split across five arrays so the shader can use scalar
layout and avoid a 128-byte struct with alignment surprises.  The
\code{nodeWritten} election buffer is device-local; it is cleared with
\code{vkCmdFillBuffer} immediately before the dispatch, not uploaded from
the host.

\subsection{Capacity management}

The input buffers grow to the largest chunk seen so far:
\begin{equation}
  \text{nodeCap}\ge |\N|,\qquad
  \text{cellCap}\ge |\text{cells}|.
\end{equation}
A grow recreates the buffers and bumps \code{bufferGeneration}.  The
descriptor writer checks the generation and rewrites bindings 0--9 only
when it changes.  In steady state, a chunk with the same capacity reuses
the descriptors and pays zero descriptor writes.  The device-local
\code{nodeWritten} buffer follows \code{nodeCap}; it is not grown per chunk.
The packed vertex/index spans belong to the renderer and are reserved per
chunk.

\subsection{Capacity plan and the provider callback}

The provider is called before the dispatch with a demand derived from the
CPU topology:
\begin{itemize}
\item vertex capacity = \code{activeVertexCount} (dense surface-node slot
  address space);
\item index capacity = $\sum_c 6\cdot\text{segments}(c)\cdot
  \text{crossEdges}(c)$;
\item the cross-edge count and maximum segment span are carried for the
  profile.
\end{itemize}
The provider reserves a stable draw-entry slot and a variable-size span:
\begin{lstlisting}[caption={Span reservation callback (shape).},label={lst:provider}]
GpuMeshSpan span = ir.reserveGpuMeshSpan(chunkId,
    demand.suggestedVertexCapacity, demand.suggestedIndexCapacity,
    &cubeMin, &cubeMax, level, &cubeMin);
target.vertexBuffer = span.vertexBuffer;
target.baseVertex = span.baseVertex;
target.vertexCapacity = span.vertexCapacity;
target.indexBuffer = span.indexBuffer;
target.firstIndex = span.firstIndex;
target.indexCapacity = span.indexCapacity;
target.slotIndex = span.slotIndex;
\end{lstlisting}
If the provider cannot fit the requested capacities, the build fails
without dispatching.  The caller releases the reservation and leaves the
revision uncommitted, so a later opportunity can retry.

\subsection{Overflow growth}

If the shader reports an index or vertex overflow, the service calls the
release callback with the reservation it just failed to use.  The renderer
frees that unpublished span immediately (no draw command references it) and
defers the free of any previously published span through the usual staged
meta flush.  The service then asks the provider for the exact reported
capacities and retries the same chunk with a freshly cleared
\code{nodeWritten} buffer.  At most three attempts are made.  Only this
chunk's spans are released and re-reserved; other chunks and the global
descriptor state are untouched.  A fixed break-list overflow or an invalid
slot fails immediately because capacity cannot fix it.

\subsection{Publication}

After the fence, the host reads the counters.  If the chunk is valid:
\begin{enumerate}
\item \code{publishGpuMeshSpan(slot, vertexExtent, indexCount)} trims the
  reserved span to the referenced extent, frees the tail, and stages the
  draw command + bounds + LoD meta;
\item the scene publisher stores a \code{LoDMesh} with the slot, counts,
  base vertex, first index, bounds, brush, and mapping metadata;
\item the renderer's indirect cull sees the entry on the next frame.
\end{enumerate}
The old span is released only after the new one is published.  The release
is staged for the cull frame after next, so no in-flight cull or draw can
reference a freed range.

\subsection{Why the host reads only counters}

The GPU writes indices and vertices.  The host reads only counters:
\begin{center}
\begin{tabular}{@{}ll@{}}
\toprule
Counter & Use\\
\midrule
\code{indexCount} & indices appended; growth requirement\\
\code{indexOverflow} & triangles rejected by the index span\\
\code{vertexExtent} & published extent; growth requirement\\
\code{vertexOverflow} & triangles rejected by the vertex span\\
\code{vertexWriteCount} & distinct payloads written; utilization\\
\code{invalidSlotCount} & flattening disagreement (fatal)\\
\code{overflowCount} & fixed-capacity recursion failure (fatal)\\
\code{droppedCount} & total rejected triangles\\
\code{findCellCalls} & diagnostics only\\
\code{reuseHits} & diagnostics only\\
\bottomrule
\end{tabular}
\end{center}
No position, normal, or 64-byte vertex ever crosses the bus.  This is the
single most important property of the design: the host cannot accidentally
stall on a vertex readback because there is no vertex readback API.

\subsection{Error policy}

A chunk build can fail in the following ways:
\begin{enumerate}
\item \textbf{flattening unavailable}: the chunk is skipped with a log;
\item \textbf{fixed-array overflow or invalid slot}: the chunk is failed and
  retried after the next edit; it is not partially drawn;
\item \textbf{span overflow}: the exact reported capacities are requested
  and the dispatch is retried once or twice, scoped to the chunk;
\item \textbf{pool span full}: release the slot and skip; the pool is sized
  for the worst case and this is a sizing bug, not a normal state.
\end{enumerate}
The revision is committed only when the chunk is actually published or
legitimately empty.  A failed build is not recorded as up-to-date; a
transient pool failure therefore retries instead of becoming a permanent
hole.
